% !TeX root = ../papers/nc-p3.tex
\section{被約二十点と有限な低階数軌跡における反対八次式の整性}
\label{sec:clean-octic-integrality}
\begin{theorem}\label{thm:clean-octic-integrality}
四辺有限の標準有限Frobenius窓で、入力の二十点は被約とする。
全階数反対平方根の階数2以下の軌跡が有限なら、その行列式は幾何学的に既約な八次式である。
従ってこの範囲の全窓に定理\ref{thm:octic-kernel-canonical-model}の標準偏極と二捩れによる復元が適用される。
\end{theorem}
\begin{proof}[接平面の族が出力を覆わない場合]
前節までの結果から、可約なら正規三次因子 $S$ があり、その接平面の一般次数 $r=m-18$ は0か1である。
$S$ の一般点 $x$ で $N(x)$ の階数三の二次錐に接する平面は、双対空間の平面 $u(x)^\perp$ 内の
滑らかな二次曲線 $B_x$ をなす。ここで $u$ は核である。$B_x$ から元の二次錐を一意に復元でき、
$x\mapsto[N(x)]$ は二次Veroneseの閉埋込みなので族の次元は2である。
また $S$ は既約三次で二次式に入らず、$[N(x)]$ の包は全九次元射影空間である。

$r=0$ なら $B_x$ の合併の閉包は整曲面 $Z$ である。平面なら全ての $u(x)$ が一致し
$L^2=27-m>0$ に反する。二次曲面なら各一般 $B_x$ はその平面切断である。
$Z$ が滑らかなら、その可逆な対称行列の逆を $Q$ とすると
\[
              [N(x)]=[(u^tQ u)Q-(Qu)(Qu)^t] .
\]
実際試験点 $v$ の双対平面 $v^\perp$ がこの切断に接する条件は
$(u^tQu)(v^tQv)-(u^tQv)^2=0$ である。
右辺の写像はVeroneseに可逆な線型写像を合成したもの（跡方向と跡零部分を分ける）で基点がない。
真変換上で両辺の超平面を比べれば $2H=2L$ となるが、$H^2=3$ と $L^2=9$ が矛盾する。
$Z$ が二次錐なら座標を選び、接触条件を
$\sum_{i=1}^3(u_4v_i-u_iv_4)^2=0$ と書ける。これは $u^\perp\cap v^\perp$ を錐の三変数平面に射影して
底の二次曲線への接触を取った式である。対応する $v$ の二次式の包は高々五次元であり、これも不可能である。

$\deg Z\ge3$ の場合、解消上の一般 $B_x$ の真変換 $B$ は滑らかな有理曲線である。
実際その滑らかな二次曲線への固有双有理射は有限なので同型である。二次元族を動くから
$h^0(\mathcal O_B(B))\ge2$（Hilbertスキームの接空間）、従って $B^2\ge1$。
超平面の引戻し $h$ に対し $hB=2$ なのでHodge指数定理から $B^2\le4/\deg Z$、従って $B^2=1$。
解消の不正則数は零である。Albanese射は有理曲線を潰し、像が曲面なら一般点はそのような繊維に属せず、
像が曲線なら繊維成分の自乗は非正であるためである。線型系 $|B|$ は完全列により三切断を持ち、
$B$ 上で全一次式を与え、基点なしである。自乗1なので $\mathbb P^2$ への双有理射となる。
族の一般元は代数的同値、従って線型同値で、その像は直線になる。
逆に $Z$ への射は平面上の共通因子のない四つの二次式で与えられる。一般直線の像が二次曲線だからである。
一般の $B_x$ に対応する直線を $a^t z=0$ とし、出力の試験点 $v$ による四つの二次式の線型結合を
$z^t Q(v)z$ と書く。接触条件は $a^t\operatorname{adj}(Q(v))a=0$ であり、一般直線上での二項二次式の判別式である。
ゆえに $N(x)$ は $\operatorname{adj}(Q(v))$ の六つの成分の包に入ってしまう。これは十次元の包に反する。
これで $r=0$ を排除した。
\end{proof}
\begin{proof}[唯一の残余点と四節点三次の二重被覆]
$r=1$ なら $n=15,b=4$ である。四節点三次は、四点を座標に取った
$\sum c_i\prod_{j\ne i}x_j=0$ の全係数が非零なので射影同値を除き唯一である。
次の六次del Pezzo曲面 $T$ からの二重被覆を使う。二次元トーラスの指標
$(1,0),(0,1),(1,1)$ とその負を頂点とする六角形の射影トーリック曲面を取る。
各頂点から出る二辺の原始ベクトルは格子基底で、六辺の格子原点からの距離は1であるから $T$ は滑らかで
超平面 $H=-K_T$、$H^2=6$。隣接する頂点は正負の対にならないので
\[
              1,\quad s+s^{-1},\quad t+t^{-1},\quad st+(st)^{-1}
\]
は境界でも基点を持たない（各トーラス軌道の面に属する単項式を残せばよい）。
これらの間の関係はアフィン座標で $x^2+y^2+z^2-xyz=4$ である。$H$ は豊富なので射は有限であり、
三次像への次数は2。トーラス反転 $\iota$ による商で、四つの固定点 $(\pm1,\pm1)$ は四節点に写る。
反転は境界の面を交換するのでほかに固定点はない。

三次核ベクトルを $T$ に引き戻す。非分岐の二十二の印は単純である。
四固定点では余接三空間が核座標の一次部で生成される。例外二次曲線がその核座標で平面を張るからである。
局所作用は滑らかな二変数の符号反転なので節点の極大イデアルの引戻しは極大イデアルの二乗となる。
従って基底長は $s_0=22+4\cdot3=34$。定理\ref{thm:low-degree-octic-spinor}の直交三束 $E$ と
その等方直線の三次元多様体 $P\to T$ を取る。$P$ 上の四切断の一般三式の完全交叉は
$2s_0-6H^2=32$ の長さを持つ。二十二の単純な印に各一、四固定点に各二の等方根があり、合計30個の相異なる基点である。
一般三式を選ぶことは出力の一般三空間を選ぶことであり、基点以外では $N|_W$ の階数1条件と同値である。
三次元の正則な二項二次式空間で、$W$ の像が共通根を持てば等方直線の直交平面に入り、逆も成立する
（$N$ が階数3のとき $W$ の像は少なくとも二次元）。
従って一般三空間では $r=1$ の唯一の非標識入力の二つの持ち上げにもそれぞれ一つの根がある。
長さ32を尽くすので完全交叉は被約で、二つの残余点は反転で交換される。

$P$ は $E$ の直線の射影束内の等方二次曲面族で、相対超平面を $\xi$ とすると四切断の束は
$D=\xi+H$ である。随伴より $K_P+D=K_T+H$ は反転同変に自明である。
実際 $\det E$ の自明化は核列の行列式から同変であり、直交形式と四つの座標切断も入力からの引戻しである。
$K_T+H$ の自明化はトーラス上の $ds/s\wedge dt/t$ と超平面切断1との積であり、
境界で極と零が打ち消し合い、反転がこれを固定する。
三式の完全交叉の標準束は $K_P+3D=2D$ の制限である。
第四の切断の平方をこの同変自明化で標準切断とすれば、周囲から延長するので完全跡は零である。
一方、基点では平方は消え、残余の二点では反転が交換するため局所跡は相等しく非零である。
三式も反転不変なので随伴で用いる順序付き定義式は保たれ、標数零でこの和は消えない。矛盾である。
残る分解がなく、平方自由性と合わせて整性を得る。
\end{proof}
